% ============================================================================
% NEW MATERIAL, October 5, 2026 (agent draft for Tommaso; not in JMP_DS_draft).
% Continues sections/02_empirical_evidence.tex, which is left unchanged.
% Register follows Boar, Gorea and Midrigan (RES 2022), Section 2: a short
% "Data" subsection that defines every variable, then facts stated one at a
% time, each tied to the model object or target it disciplines.
% Every number carries its source in a comment: [DATA] = measured moment,
% with builder; [MODEL] = model output, with run.
% ============================================================================
\providecommand{\mocktodo}[1]{\textcolor{red}{[TODO: #1]}}

\subsection{Fertility: level, timing and income}
\label{sec:emp-fertility}

The facts above describe how housing adjusts once a child arrives. The model
also has to reproduce when children arrive and who has them. This subsection
documents three facts about fertility itself: completed fertility is close to
two children with one woman in five childless, first births are concentrated
in the mid twenties, and early childbearing is concentrated in lower-income
families, while completed fertility is nearly flat in income.

\subsubsection*{Data.}
Fertility levels and timing come from two sources. The June Fertility
Supplement of the Current Population Survey (CPS) records the number of
children ever born to each woman. I use the 2004 and 2006
supplements for the calibration targets, so that they refer to the same years
as the wealth and housing targets, and the 2024 supplement for the income
gradient, because it is the most recent.\mocktodo{the gradient and the
targets come from different CPS years; recompute the gradient on 2004 and 2006
if the paper keeps both.} Completed
fertility refers to women aged 40 to 44. The timing of first births comes from
the National Center for Health Statistics (NCHS) natality files, which record
every live birth in the United States with the mother's age and the birth
order; I pool first births in 2003 to 2006. Family income in the CPS is the
combined money income of the family over the previous twelve months, reported
in brackets. All CPS statistics use the person weights and cap children ever
born at three, the largest number the model distinguishes.%
\footnote{Capping matters for levels: among women aged 40 to 44 in the 2024
CPS, capping at five instead of three raises mean children ever born by 0.13
to 0.23 across income groups. \mocktodo{add the uncapped values and the CPS builders to the data
appendix: code/data/cps\_fertility/build\_early\_fertility\_target.py,
output/model/credit\_mechanism\_20261004/measurement/recompute.py.}}

\subsubsection*{Completed fertility and childlessness.}
Among women aged 40 to 44 in the 2004 and 2006 CPS, 19.8 percent have no
children, and 21.4 percent of mothers have exactly one.
% [DATA] cps_childlessness 0.19828, cps_exactly_one 0.21366; CPS June 2004/06,
% women 40-44, FREVER 0-20, supplement weight. Chain 11 target rows 2-3.
In the 2024 supplement, mean children ever born is 1.92 and 18.8 percent are
childless.
% [DATA] CPS June 2024, women 40-44, n = 3,278; 1.9184 (SE 0.027), 0.1882
% (SE 0.008); code/data/cps_fertility/build_cps_fertility_targets.py.
The two margins matter separately in the model: childlessness is the extensive
margin, governed by the first birth, and the share with exactly one child is
the first intensive margin, governed by whether parents continue.

\subsubsection*{Timing.}
First births are concentrated in the mid twenties. The mean age of mothers at
first birth is 26.0 years in the 2003--2006 natality files, and 24.9 percent
of first births occur at age 30 or later.
% [DATA] nchs_mean_age 25.976 (SE 0.15 declared), nchs_share30 0.249;
% code/data/nchs_natality_timing/build_first_birth_timing.R; age cells mapped
% to four-year midpoints. Chain 11 rows 4 (target) and 5 (validation).
A woman aged 25 in the 2004 and 2006 CPS has had 0.81 children on average.
% [DATA] early_fertility 0.80953, SE 0.028, n = 1,774; CPS June 2004/06,
% women at exact age 25, capped at 3. Chain 11 row 14.
The age at which childbearing starts matters here because it coincides with
the ages at which households form, accumulate their first savings, and buy
their first home.\mocktodo{a sentence comparing the wealth of childless
renters at 25 and at 32 (PSID) would make this concrete; not yet measured.}

\subsubsection*{Income.}
Early childbearing is concentrated in lower-income families.
Table~\ref{tab:emp-fertility-income} groups women by family income into thirds.
Among women aged 24 to 26, those in the bottom third have had 0.60 children,
those in the top third 0.23. Among women aged 40 to 44 the gradient has
nearly disappeared: 1.78 children in the bottom third and 1.73 in the top,
a difference of $-0.05$ with a standard error of $0.05$. Lower-income women
therefore do not have many more children over their lives; they have them
earlier.

\begin{table}[t]
\centering
\caption{Children ever born by family income (CPS, 2024)}
\label{tab:emp-fertility-income}
\begin{tabular}{lcccc}
\toprule
 & Bottom third & Middle third & Top third & Top minus bottom \\
\midrule
Women aged 24--26 & 0.596 & 0.448 & 0.226 & $-0.369$ \\
 & & & & $(0.052)$ \\
Women aged 40--44 & 1.784 & 1.702 & 1.733 & $-0.051$ \\
 & & & & $(0.050)$ \\
\addlinespace
\multicolumn{5}{l}{\emph{Reference person or spouse only}}\\
Women aged 24--26 & 1.214 & 0.923 & 0.569 & $-0.645$ \\
 & & & & $(0.096)$ \\
Women aged 40--44 & 2.093 & 1.911 & 1.875 & $-0.217$ \\
 & & & & $(0.048)$ \\
\bottomrule
\end{tabular}
\par\smallskip
\begin{minipage}{0.92\linewidth}
\footnotesize\textit{Notes:} June 2024 CPS Fertility Supplement. Mean number
of children ever born, capped at three, weighted by the CPS person weight.
Thirds are equal-weight groups of family money income in the previous twelve
months; women in an income bracket that straddles a cut-off are split across
groups in proportion to the bracket's mass. Standard errors in parentheses are
approximate, clustered by household, and ignore the CPS sample design. 1,673
women aged 24--26 and 3,278 aged 40--44.
% [DATA] output/model/credit_mechanism_20261004/measurement/tables.json,
% keys age24_26_all_women_fractional_weighted_thirds etc.; AUDIT.md there.
% Status: "cached_only_diagnostic_not_calibration". Not a calibration target.
\end{minipage}
\end{table}

Two features of the measurement limit how far this fact can be pushed. First,
a woman of 25 who lives with her parents reports her parents' family income,
which overstates the resources of young childless women. Restricting to women
who head a household or are married to its head, as in the second panel,
raises fertility at every income level and makes the young gradient steeper,
not flatter. Second, a birth lowers a mother's earnings, so part of the
gradient at young ages runs from children to income. The fact is therefore a
description of who has children early, not an estimate of how income affects
fertility, and I use it to evaluate the model rather than to calibrate it.
% NOTE for Tommaso: chain 13 gives 0.124/0.508/0.973 (24-26) and
% 1.318/1.782/2.104 (40-44) by thirds of the persistent earnings state
% [MODEL, chain 13, fixed price; thread_credit_null/README.md]. The model
% ranks by current earnings state z, the CPS by family income. This is the
% largest untargeted miss; it belongs in the quantification, not here.

\subsubsection*{Ownership around parenthood.}
The PSID event study above shows ownership rising by 11 percentage points in
the two years around the first birth and by 16 points three to four years
later, on a base of 41 percent. The cross-section shows the same gap among
established households. In the 2005--2006 American Community Survey (ACS),
household heads aged 30 to 55 whose oldest child is under four are 12.8
percentage points more likely to own than heads with no children, and 67.6
percent of all heads in this age range own their home.
% [DATA] recent_parent_ownership 0.12761, ownership_30_55 0.67626; ACS
% 2005-06, HHWT; code/empirical/housing/build_initial_housing_profile_diagnostic.py.
% Chain 11 rows 10 and 13.
\mocktodo{state the dwelling-structure restriction used in both ACS ownership
moments (UNITSSTR codes 3--10 in the builder) in words, and say whether it is
also imposed on the model observer.}

\subsubsection*{Facts used in the quantification.}
Table~\ref{tab:emp-targets} collects the measured moments the model is asked
to match, with their sources and samples. The quantification
explains which parameter each one disciplines. Two of them are new relative to
the housing evidence above: the number of children by age 25, which
summarizes how early families form, and the room response to the first birth,
measured in the sample that fixes household position before the birth.

\begin{table}[t]
\centering
\caption{Measured moments}
\label{tab:emp-targets}
\small
\begin{tabularx}{\linewidth}{Xccl}
\toprule
Moment & Value & Data & Sample \\
\midrule
\multicolumn{4}{l}{\emph{Fertility}}\\
Childless, women 40--44 & 0.198 & CPS 2004, 2006 & women 40--44 \\
Exactly one child, among mothers 40--44 & 0.214 & CPS 2004, 2006 & women 40--44 \\
Children ever born at age 25 & 0.810 & CPS 2004, 2006 & women aged 25 \\
Mean age at first birth & 25.98 & NCHS 2003--2006 & all first births \\
Share of first births at 30 or later$^{\dagger}$ & 0.249 & NCHS 2003--2006 & all first births \\
\addlinespace
\multicolumn{4}{l}{\emph{Housing}}\\
Mean rooms & 5.73 & AHS 2007 & heads 18--85 \\
Ownership rate & 0.676 & ACS 2005--2006 & heads 30--55 \\
Ownership gap, recent parents minus childless & 0.128 & ACS 2005--2006 & heads 30--55 \\
Rooms added three to four years after first birth & 1.465 & PSID 1968--2019 & see notes \\
Rooms, three or more children minus one or two$^{\dagger}$ & 0.385 & ACS 2005--2006 & heads 30--55 \\
\addlinespace
\multicolumn{4}{l}{\emph{Wealth}}\\
Net worth over gross labor earnings & 4.46 & PSID 2005, 2007 & heads 18--85 \\
Bequest flow over net worth & 0.0073 & SCF 2007 & all households \\
Wealth-to-income, p90 over median, ages 76--84$^{\dagger}$ & 3.52 & PSID 2003--2005 & heads 76--84 \\
\bottomrule
\end{tabularx}
\par\smallskip
\begin{minipage}{0.95\linewidth}
\footnotesize\textit{Notes:} $^{\dagger}$ Not targeted; used to evaluate the
model. Children ever born are capped at three. The room response is the
Sun--Abraham event-study estimate for calendar years three and four after the
first biological birth relative to years two and three before it, in the
sample of adults who were reference person or spouse two to three years before
the birth (117,853 person-years; standard error 0.050); the household sample
of Figure~\ref{fig:emp-rooms-path} gives 1.03. Net worth excludes business and
farm equity and other real estate. The bequest flow is the estate flow
directed to children.
\mocktodo{confirm the wealth numerator and the 1.465 sample with Tommaso before
this table leaves the mock; both changed in the last two weeks.}
% [DATA] all values are the chain 11 target contract (target fingerprint
% c7a3d185...), from CALIBRATION_STATUS.md and
% output/model/overnight_dual_20261004_v1/terminal_collection/a_provisional/
% target_fit_new_contract.csv. Builders: AHS code/data/ahs_supply_snapshot/
% build_ahs_2007_room_target.py; PSID rooms code/data/psid_followup_mar2026/
% sa_rooms_first_birth_v2.do (arm A2h); SCF code/data/scf/build_bequest_flow_2007.py.
\end{minipage}
\end{table}

\subsubsection*{Summary.}
Children arrive early, mostly in the mid twenties, and earliest in
lower-income families. Their arrival coincides with a move into a larger
dwelling and, for many families, into ownership. The model in the next
section is built so that both decisions are made by the same household, at the
same ages, out of the same limited savings.
